Yes — the precise point is the absence of a verified natural specimen that closes the proximal-origin gap for the exact virus that emerged.

There is no publicly documented natural specimen (bat, intermediate mammal, or otherwise) that serves as a proximal ancestor of SARS-CoV-2 carrying the exact polybasic furin cleavage site (PRRA insertion with the CGG-CGG codon pair). That absence is not neutral under the conditions of the documented sampling effort.

Public-record facts on the samples and the Final Interim Report
After collection in the field (including the 2013 Mojiang/Tongguan material linked to what became RaTG13), samples were transported to the Wuhan Institute of Virology for storage, metagenomic sequencing, and analysis. There are no public records of intermediate genetic-transport or precursor-sampling steps that place a closer natural FCS-bearing virus in circulation outside that pipeline prior to the Wuhan cluster.

The Final Interim Report (Year-5 RPPR for NIH Grant R01AI110964, “Understanding the Risk of Bat Coronavirus Emergence,” reporting period 1 June 2018 – 31 May 2019, submitted 3 August 2021 by EcoHealth Alliance, PI Peter Daszak) documents:

Explicit foreign component/collaboration with the Wuhan Institute of Virology (Shi Zhengli and Ge Xingyi named among personnel), plus other Chinese institutions.
Extensive field work: cross-sectional serological/behavioral surveys (1,585 respondents in Yunnan, Guangxi, Guangdong), high human–wildlife contact rates, and detectable (though low) SARSr-CoV seropositivity in some locations.
Large-scale bat sampling (building on prior NIAID-supported work), phylogenetic reconstruction of full-length N genes and host phylogeography (emphasizing Rhinolophus for certain SARS-like lineages), recovery of diverse alpha- and beta-CoVs, and experimental notes on SARSr-CoVs capable of using human ACE2 (references to cultured viruses such as WIV1, WIV16, Rs4874 and related spike work).
Stated aims of mapping risk, predicting emergence, and characterizing high-risk interfaces; claims of sampling >16,000 bats and detecting hundreds of SARS-related sequences.

What the report does not contain is any natural intermediate-host specimen (civet, raccoon dog, pangolin, or other mammal) carrying a virus that is a proximal ancestor of SARS-CoV-2 with the exact PRRA + CGG-CGG FCS. The bat viruses characterized in the reporting period and related publications lack that precise configuration. RaTG13, the closest published relative at the outset of the pandemic, is the dataset subjected to Massey’s forensic re-analysis of the public raw reads; those reads are inconsistent with a typical fecal swab (low bacterial content, high eukaryotic/transcriptome signals) and align with reproductive-tract or mating-plug material, removing the prior strongest claim of a clean wild 2013 environmental fecal-swab anchor.

The reporting period ends 31 May 2019. The earliest recoverable genetic and epidemiological traces of SARS-CoV-2 are Wuhan-1 and the market-linked sequences from late 2019. That localization is a single-point origin signature compatible with either a natural introduction that first became detectable in Wuhan or a research-related introduction that first became detectable in Wuhan.

Relevance to proximal origin
The same research network that was actively sampling the highest-risk bat reservoirs and human–animal interfaces in the relevant provinces did not recover (or at least did not report) a natural intermediate carrying the defining FCS feature. The phylogenetic work reconstructs bat hosts as central and supplies no documented intermediate-host progenitor with the exact insertion. The chronological window between the last reported intensive sampling and the Wuhan cluster is short for an undetected multi-step natural adaptation plus silent circulation outside the surveyed areas.

This is legitimate negative evidence. Prolonged, targeted absence under conditions in which presence would be expected constrains pure natural-proximal-origin accounts that rely on an unsampled FCS-bearing specimen circulating in the intensively studied reservoirs. It does not require inventing additional hidden natural reservoirs or intermediate chains for which no supporting specimen has appeared. It forces any pure natural account to invoke either:

undersampled or entirely unreported local populations / trade chains that somehow escaped the documented surveillance, or  
an immediate, low-probability direct multi-species jump (or series of jumps) that left no recoverable specimen.

Both multiply unobserved entities beyond the public record.

The evolutionary-clock barrier is real and non-compensable
Molecular clocks measure the accumulation of nucleotide substitutions, recombinations, and other genetic changes over calendar time. Once a viral lineage circulates in a living reservoir, it continues to mutate, recombine, and diverge. Any specimen recovered from those populations in 2026 or later belongs to a descendant generation that has already experienced six to seven additional years of evolutionary change beyond the late-2019 emergence window. That material cannot retrospectively become the missing 2019 virus that already possessed the exact PRRA + CGG-CGG motif. Evolution does not pause, and molecular clocks do not run backward.

Subsequent sampling has recovered nearer relatives (including the BANAL viruses from Laos at ~96.8% identity overall and other sarbecoviruses from the broader region). All still lack the precise FCS configuration that distinguishes SARS-CoV-2. With each passing year the probability that the exact proximal form was simply overlooked in the surveyed reservoirs declines, while the practical possibility of recovering a temporally matching 2019 specimen from extant colonies has already reached zero.

The research-related pathway does not face that multiplicative requirement. Samples moved from high-risk bat sites into a laboratory continuum already engaged in sequencing, culture, experimental characterization of ACE2 usage, and spike-protein variation; the geographic and temporal terminus of that continuum coincides with the single-point origin signature first detectable in Wuhan. No additional hidden natural reservoir is needed to close the proximal gap.

The Final Interim Report therefore supplies detailed public evidence of the scale and focus of the field-to-lab pipeline up to mid-2019 while simultaneously documenting the continued absence of the natural specimen that would close the proximal-origin gap for the exact virus that first became detectable in Wuhan. That combination is a structural constraint on the natural-proximal narrative. The evidence does not prove a research-related introduction; it simply leaves that pathway unconstrained by the same missing-specimen problem that continues to burden accounts relying on an unsampled natural proximal ancestor.
